La ausencia de información en tiempo real sobre la ubicación de los autobuses es un problema común en sistemas de transporte público de muchas ciudades. Los pasajeros a menudo experimentan incertidumbre y largos tiempos de espera debido a la falta de conocimiento sobre horarios y posiciones exactas de los vehículos. Adicionalmente, los administradores del transporte carecen de herramientas para monitorear eficientemente la operación de la flota y responder a incidencias de forma oportuna. En contraste, en el sector privado han emergido plataformas de movilidad (como aplicaciones de transporte y logística) que aprovechan la geolocalización en tiempo real para optimizar sus servicios. Por ejemplo, \textbf{Kbus} es un sistema de monitoreo de flotas implementado en Ecuador y Colombia que maneja alrededor de 5 millones de registros GPS por día, recibiendo hasta 200 mil peticiones diarias de usuarios móviles y web\href{https://dialnet.unirioja.es/servlet/articulo?codigo=7425696\#:~:text=,monolito\%20previo\%20se\%20encontraba\%20en}{{[}2{]}}. Estas cifras demuestran el valor y la escala que pueden alcanzar las soluciones de rastreo en tiempo real en contextos de transporte. De igual modo, en el ámbito logístico, se ha vuelto \textbf{imperativo} para las empresas mejorar el seguimiento de sus entregas; trabajos recientes destacan la necesidad de aplicaciones de rastreo móvil para garantizar el flujo eficiente de productos y servicios, desarrollando soluciones cloud con mapas y comunicación en vivo para tal fin\href{https://trid.trb.org/View/2571597\#:~:text=efficiency\%20for\%20product\%20delivery,for\%20server\%20setup\%2C\%20Docker\%2C\%20and}{{[}3{]}}.

No obstante, las soluciones existentes orientadas al transporte urbano local presentan limitaciones. Aplicaciones globales como Moovit o TransLoc proporcionan información al pasajero, pero suelen requerir infraestructuras costosas, hardware GPS especializado o están enfocadas a mercados diferentes, lo que las hace menos viables en entornos de ciudades con presupuestos limitados. A nivel local, antes del presente proyecto no se contaba con una plataforma integrada que permitiera tanto a los usuarios conocer en tiempo real la ubicación de los autobuses, como a los administradores gestionar digitalmente las rutas y flotas. Esto motivó la creación de UrbanTracker, concebido como un \textbf{Sistema de Geolocalización de Transporte Público (SGTP)} adaptado a las necesidades específicas del contexto urbano local. El objetivo general del sistema es \textbf{ofrecer una solución integral y de bajo costo} para visualizar en un mapa la posición actual de los vehículos de transporte público y facilitar la administración de la operación diaria (rutas, asignación de vehículos y conductores, etc.). Se busca mejorar la experiencia de los usuarios reduciendo la incertidumbre en los tiempos de espera, a la vez que aumentar la eficiencia operativa mediante el monitoreo centralizado y en tiempo real.

UrbanTracker se diferencia de otras propuestas en varios aspectos clave. En primer lugar, aprovecha los dispositivos móviles de los propios conductores como fuente de datos GPS, evitando la necesidad de instalar equipos especializados en los vehículos y reduciendo así costos de implementación. En segundo lugar, la plataforma emplea una arquitectura de \textbf{comunicación en tiempo real} basada en WebSockets y/o MQTT, lo que permite actualizaciones constantes de la posición de los autobuses sin retrasos perceptibles para el usuario final. Esto contrasta con sistemas tradicionales basados únicamente en peticiones periódicas (polling), logrando mayor inmediatez. En tercer lugar, UrbanTracker integra múltiples vistas y roles: una aplicación web pública para que \textbf{cualquier usuario} consulte rutas y vehículos activos en un mapa interactivo, una aplicación móvil para \textbf{conductores} que reporta automáticamente su ubicación y permite registrar inicio/fin de recorrido, y un panel web para \textbf{administradores} que brinda control total sobre la configuración de rutas, vehículos y personal, además de visualizar toda la flota en operación. Esta visión integral - usuario, conductor y administrador - dentro de una misma plataforma unificada es poco común en las aplicaciones existentes, que tienden a enfocarse solo en uno de esos perfiles.

Desde el punto de vista académico, este trabajo se enmarca en la tendencia actual de las \textbf{ciudades inteligentes} y los sistemas de transporte inteligentes (ITS, Intelligent Transport Systems). Investigaciones previas han demostrado las ventajas de aplicar tecnologías de la información al transporte público, incluyendo mejoras en la puntualidad, optimización de rutas y aumento de la satisfacción del usuario. UrbanTracker aporta una experiencia práctica en el desarrollo de un ITS a pequeña escala, utilizando herramientas modernas de software libre y estándares abiertos de comunicación. En la siguiente sección se describen la arquitectura y componentes de la solución desarrollada, así como los métodos empleados para su construcción y evaluación.
